% Part: sets-functions-relations
% Chapter: infinite
% Section: dedekind
%
\documentclass[../../../include/open-logic-section]{subfiles}

\begin{document}

\olfileid[fr]{sfr}{infinite}{dedekind}
\olsection{Algèbres de Dedekind}

Nous ne voulons pas seulement que les entiers naturels soient en nombre
infini ; nous voulons aussi qu'ils possèdent certaines propriétés
algébriques : l'addition, la multiplication et les autres opérations
doivent se comporter comme il convient.

L'idée de Dedekind consiste à prendre pour notion fondamentale la
\emph{fonction successeur}, puis à caractériser les nombres par les
propriétés suivantes :
\begin{enumerate}
	\item Il existe un nombre, $0$, qui n'est le successeur d'aucun nombre,
	\\c'est-à-dire $0 \notin \ran{s}$,
	\\c'est-à-dire $\forall x\ s(x) \neq 0$.
	\item Des nombres distincts ont des successeurs distincts,
	\\c'est-à-dire que $s$ est !!a{injection},
	\\c'est-à-dire $\forall x \forall y (s(x) = s(y) \lif x = y)$.
	\item\ollabel{repeatedapplication} Tout nombre s'obtient à partir de
	$0$ par applications répétées de la fonction successeur.
\end{enumerate}
Les deux premières conditions s'expriment aisément en logique du premier
ordre, comme on vient de le voir. Mais la condition
\olref{repeatedapplication} ne peut pas s'exprimer à l'aide de la seule
logique du premier ordre dans ce langage arithmétique. L'avancée décisive
de Dedekind consiste à reformuler la condition
\olref{repeatedapplication} en termes ensemblistes, de la manière suivante :
\begin{enumerate}
	\item[3$'$.] Les entiers naturels forment le plus petit ensemble
	contenant $0$ et \emph{stable par la fonction successeur} :
	appliquer $s$ à un !!{element} de cet ensemble donne encore
	un !!{element} de l'ensemble.
\end{enumerate}
Il nous faut toutefois expliciter cette idée pas à pas.\footnote{Note de
l'édition française : la condition « contenant $0$ », omise dans cette
reformulation de la source, est nécessaire ; l'ensemble vide est lui aussi
stable par le successeur. La reformulation ensembliste peut être écrite
dans le langage du premier ordre de la théorie des ensembles, qui permet
de quantifier sur les ensembles.}

\begin{defn}\ollabel{Closure}
	Soit $f\colon A\to A$. Une partie $X\subseteq A$ est
	\emph{stable par $f$} si et seulement si
	$(\forall x \in X)f(x) \in X$. Pour tout $o\in A$, on définit la
	\emph{clôture de $o$ sous $f$} par
	$$\closureofunder{f}{o} = \bigcap\Setabs{X\subseteq A}{o \in X\text{ et }X
\text{ est stable par }f}.$$
\end{defn}
\noindent\emph{Précision de l'édition française.} Nous fixons ici l'ensemble
ambiant $A$ et supposons que $f$ est une application de $A$ dans lui-même.
Ces hypothèses, laissées implicites dans la source, garantissent que la
famille dont on prend l'intersection est un ensemble non vide.

Ainsi, $\closureofunder{f}{o}$ est l'intersection de toutes les parties
de $A$ stables par $f$ qui contiennent $o$ comme !!{element}.
Intuitivement, $\closureofunder{f}{o}$ est donc la \emph{plus petite}
partie de $A$ stable par $f$ qui contient $o$ comme !!{element}.
Le résultat suivant donne un sens précis à cette intuition.
\begin{lem}\ollabel{closureproperties}
	Pour toute application $f\colon A\to A$ et tout $o \in A$ :
	\begin{enumerate}
		\item\ollabel{closurehaselem} $o \in \closureofunder{f}{o}$ ;
		\item\ollabel{closureclosed} $\closureofunder{f}{o}$ est stable par $f$ ;
		\item\ollabel{closuresmallest} si $X\subseteq A$ est stable par $f$
		et si $o \in X$, alors $\closureofunder{f}{o} \subseteq X$.
	\end{enumerate}
\end{lem}

\begin{proof}
Il existe au moins une partie de $A$ stable par $f$ qui contient $o$
comme !!{element}, à savoir $\ran{f}\cup\{o\}$.
En effet, $f$ envoie chacun de ses éléments dans $\ran{f}$.
L'intersection $\closureofunder{f}{o}$ de \emph{toutes} ces parties
existe donc. Il reste à vérifier les points
\olref{closurehaselem}--\olref{closuresmallest}.

Pour le point \olref{closurehaselem} :
$o \in \closureofunder{f}{o}$, car il s'agit de l'intersection
d'ensembles qui contiennent tous $o$ comme !!{element}.

Pour le point \olref{closureclosed}, supposons que
$x \in \closureofunder{f}{o}$. Si $X\subseteq A$ contient $o$
et est stable par $f$, alors $x \in X$ ; on a donc aussi $f(x) \in X$,
puisque $X$ est stable par $f$. Ainsi,
$f(x) \in \closureofunder{f}{o}$.

Pour le point \olref{closuresmallest}, on utilise le fait général
que, si $X \in C$, alors $\bigcap C \subseteq X$.
\end{proof}

Nous pouvons maintenant donner la définition suivante.

\begin{defn}
Une \emph{algèbre de Dedekind} est un ensemble $A$ muni d'une application
$f\colon A\to A$ et d'un élément $o\in A$ tels que :
	\begin{enumerate}
		\item \ollabel{ded:proper} $o \notin \ran{f}$ ;
		\item \ollabel{ded:injection} $f$ est !!a{injection} ;
		\item \ollabel{ded:closure} $A = \closureofunder{f}{o}$.
	\end{enumerate}
\end{defn}

Puisque $A = \closureofunder{f}{o}$, le résultat précédent montre que
$A$ est la plus petite partie stable par $f$ qui contient $o$ comme
!!{element}. Une algèbre de Dedekind est manifestement infinie au sens
de Dedekind : il suffit de regarder les conditions \olref{ded:proper}
et \olref{ded:injection} de la définition. Mais le fait le plus
remarquable est que tout ensemble infini au sens de Dedekind permet de
construire une algèbre de Dedekind.\footnote{Note de l'édition française :
la construction porte sur une partie de l'ensemble infini donné.
La formulation de la source pouvait laisser entendre que l'ensemble
tout entier recevait cette structure ; la preuve établit l'énoncé
d'existence ci-dessous en extrayant la clôture d'un élément.}

\begin{thm}\ollabel{thm:DedekindInfiniteAlgebra}
S'il existe un ensemble infini au sens de Dedekind, alors il existe une
algèbre de Dedekind.
\end{thm}

\begin{proof}
Soit $D$ un ensemble infini au sens de Dedekind. Il existe donc une
injection $g\colon D\to D$ et un élément $o\in D\setminus\ran{g}$.
Posons $A = \closureofunder{g}{o}$ ; d'après le
\olref{closureproperties}, $A$ existe et $o \in A$.
Posons $f = \funrestrictionto{g}{A}$. Comme $A$ est stable par $g$,
$f$ est une application de $A$ dans $A$. Montrons que le triplet
$A,f,o$ est une algèbre de Dedekind.

Pour la condition \olref{ded:proper}, on a $o \notin \ran{g}$ et
$\ran{f} \subseteq \ran{g}$ ; par conséquent, $o\notin\ran{f}$.

Pour la condition \olref{ded:injection}, $g$ est une injection sur $D$ ;
sa restriction $f \subseteq g$ est donc elle aussi injective.

Pour la condition \olref{ded:closure}, le
\olref{closureproperties} assure que $A$ est stable par $g$, donc
aussi par $f$. Ainsi, $\closureofunder{f}{o}\subseteq A$, d'après le
\olref{closureproperties}. De plus, $\closureofunder{f}{o}$ est stable
par $f$ ; comme c'est une partie de $A$ et que
$f = \funrestrictionto{g}{A}$, elle est également stable par $g$.
Puisqu'elle contient $o$, le \olref{closureproperties} donne
$A\subseteq\closureofunder{f}{o}$. Les deux inclusions établissent
l'égalité voulue.
\end{proof}

\end{document}
